\documentclass[12pt,a4paper]{article}

% ==================================================
% 1. CONFIGURATION
% ==================================================

% Enable / disable title page
\newif\ifshowtitlepage
\showtitlepagetrue
% \showtitlepagefalse

% Document information
\newcommand{\worktitle}{Sprint 1 Progress Report}
\newcommand{\worksubtitle}{NAO Robot Exercise System}
\newcommand{\teamname}{NAOMI}
\newcommand{\subject}{NAO Robot Exercise System}

\newcommand{\university}{Slovak University of Technology in Bratislava}
\newcommand{\faculty}{Faculty of Informatics and Information Technologies}

\newcommand{\academicyear}{2026/2027}
\newcommand{\location}{Bratislava}
\newcommand{\documentyear}{2026/2027}

% Team leaders
\newcommand{\leaderone}{Ing. Jakub Perdek, PhD.}
\newcommand{\leadertwo}{doc. Ing. Fedor Lehocki, PhD.}

% Team members
\newcommand{\memberone}{Bc. Akhmadinurov Amal}
\newcommand{\membertwo}{Bc. Danilov Denis}
\newcommand{\memberthree}{Bc. Dmitriev Anton}
\newcommand{\memberfour}{Bc. Gainetdinov Ruslan}
\newcommand{\memberfive}{Bc. Kuzin Timofei}
\newcommand{\membersix}{Bc. Tikhoblazhenko Alena}
\newcommand{\memberseven}{Bc. Zimnitski Aliaksei}

% Sprint information -- FILL IN
\newcommand{\sprintname}{Sprint 1}
\newcommand{\sprintperiod}{26 October -- 9 November 2026}
\newcommand{\statusdate}{9 November 2026}  % sprint end; statuses refer to this date


% ==================================================
% 2. PACKAGES
% ==================================================

\usepackage[
    a4paper,
    top=3cm,
    bottom=2.7cm,
    left=2.8cm,
    right=2.8cm,
    headheight=45pt,
    headsep=20pt
]{geometry}

\usepackage[T1]{fontenc}
\usepackage[utf8]{inputenc}
\usepackage[english]{babel}
\usepackage{lmodern}

\usepackage{graphicx}
\graphicspath{{../../assets/}{assets/}}
\usepackage[table]{xcolor}
\usepackage{fancyhdr}
\usepackage{titlesec}
\usepackage{setspace}
\usepackage{lastpage}
\usepackage{array}
\usepackage{float}
\usepackage{caption}
\usepackage{booktabs}
\usepackage{tabularx}
\usepackage{csquotes}

% No bibliography in progress reports, so biblatex is not loaded.

\usepackage[hidelinks]{hyperref}


% ==================================================
% 3. COLORS
% ==================================================

\definecolor{primary}{HTML}{000000}
\definecolor{secondary}{HTML}{000000}
\definecolor{muted}{HTML}{000000}
\definecolor{linecolor}{HTML}{000000}
\definecolor{naomiTableHeader}{HTML}{17445A}
\definecolor{naomiTableStripe}{HTML}{EDF7FA}
\definecolor{naomiTableRule}{HTML}{66BBD0}

\newcolumntype{L}[1]{>{\raggedright\arraybackslash}p{#1}}
\newcolumntype{Y}{>{\raggedright\arraybackslash}X}


% ==================================================
% 4. TYPOGRAPHY
% ==================================================

\setstretch{1.2}
\setlength{\parindent}{0pt}
\setlength{\parskip}{7pt}

\titleformat{\section}{\Large\bfseries\color{primary}}{\thesection}{0.8em}{}
\titleformat{\subsection}{\large\bfseries\color{primary}}{\thesubsection}{0.8em}{}
\titleformat{\subsubsection}{\normalsize\bfseries\color{primary}}{\thesubsubsection}{0.8em}{}

\titlespacing*{\section}{0pt}{20pt}{10pt}
\titlespacing*{\subsection}{0pt}{14pt}{7pt}
\titlespacing*{\subsubsection}{0pt}{12pt}{5pt}

\captionsetup{font=small, labelfont=bf, justification=centering}


% ==================================================
% 5. HEADER AND FOOTER
% ==================================================

\newcommand{\naomiheaderfooter}{%
    \fancyhf{}
    \fancyhead[L]{\includegraphics[width=4cm,height=1.15cm,keepaspectratio]{horizontal.png}}
    \fancyhead[R]{%
        \begin{tabular}{r}
            {\footnotesize\bfseries \worksubtitle}\\[3pt]
            {\scriptsize Academic Year \academicyear}
        \end{tabular}}
    \fancyfoot[L]{{\scriptsize \teamname\enspace|\enspace\subject}}
    \fancyfoot[R]{{\scriptsize \thepage\enspace/\enspace\pageref{LastPage}}}
    \renewcommand{\headrulewidth}{0.7pt}
    \renewcommand{\footrulewidth}{0.4pt}%
}

\pagestyle{fancy}
\naomiheaderfooter
\fancypagestyle{plain}{\naomiheaderfooter}


% ==================================================
% 6. TITLE PAGE
% ==================================================

\newcommand{\makemytitlepage}{
    \begin{titlepage}
        \thispagestyle{empty}
        \centering

        % University logo
        \vspace*{0.3cm}
        \includegraphics[width=6cm,height=5cm,keepaspectratio]{vertical.png}

        \vspace{0.5cm}
        {\large\bfseries \university\par}
        \vspace{0.2cm}
        {\normalsize \faculty\par}

        % Project and document title
        \vspace{1.7cm}
        {\small\MakeUppercase{\worksubtitle}\par}
        \vspace{0.5cm}
        \begin{minipage}{0.95\textwidth}
            \centering
            {\fontsize{25}{32}\selectfont\bfseries \worktitle\par}
        \end{minipage}

        % Sprint period and author
        \vspace{0.6cm}
        {\normalsize \sprintperiod\par}

        % Team information
        \vfill
        \begin{minipage}{0.9\textwidth}
            \centering
            \begin{minipage}[t]{0.46\linewidth}
                \centering
                {\normalsize\bfseries Team Leaders\par}
                \vspace{0.25cm}
                {\small
                \leaderone\par
                \leadertwo\par}
            \end{minipage}%
            \hfill
            \begin{minipage}[t]{0.48\linewidth}
                \centering
                {\normalsize\bfseries Team Members\par}
                \vspace{0.25cm}
                {\small
                \memberone\par
                \membertwo\par
                \memberthree\par
                \memberfour\par
                \memberfive\par
                \membersix\par
                \memberseven\par}
            \end{minipage}
        \end{minipage}

        % Bottom information
        \vfill
        {\small \location\enspace\textbullet\enspace\academicyear\par}
    \end{titlepage}
}


% ==================================================
% 7. IMAGES
% ==================================================

% \imageplaceholder{Text}{Caption} -- until the screenshot is added
\newcommand{\imageplaceholder}[2]{
    \begin{figure}[H]
        \centering
        \begingroup
        \setlength{\fboxsep}{0pt}
        \fbox{
            \begin{minipage}[c][5cm][c]{0.85\textwidth}
                \centering
                {\Large\bfseries IMAGE}
                \vspace{0.3cm}

                {\normalsize #1}
            \end{minipage}
        }
        \endgroup
        \caption{#2}
    \end{figure}
}

% \documentimage{path}{caption}{label}
\newcommand{\documentimage}[3]{
    \begin{figure}[H]
        \centering
        \includegraphics[width=0.85\textwidth,keepaspectratio]{#1}
        \caption{#2}
        \label{#3}
    \end{figure}
}


% ==================================================
% 8. TABLES
% ==================================================

\newenvironment{naomitable}[2]{%
    \begin{table}[H]
        \centering
        \caption{#1}
        \label{#2}
        \begingroup
        \small
        \renewcommand{\arraystretch}{1.4}
        \setlength{\tabcolsep}{6pt}
        \rowcolors{2}{naomiTableStripe}{white}
}{%
        \endgroup
    \end{table}
}


% ==================================================
% 9. DOCUMENT
% ==================================================

\begin{document}

% Title page; without it, a short heading is shown instead
\ifshowtitlepage
    \makemytitlepage
\fi


% ==================================================
% SUMMARY
% ==================================================

\section{Summary}

During \sprintname{} (\sprintperiod), the \teamname{} team completed
\textbf{12 of 14 tasks} with \textbf{36 of 44 story points}. The two
remaining tasks (SCRUM-55 and SCRUM-58, 8 story points) were finished in
code, but their pull requests were approved only on the day after the sprint
ended, so they moved to Sprint~2 and were closed there.

After the preparatory Sprint~0, \sprintname{} focused on the existing
solution. In the meetings with the supervisor, the team agreed to refactor
and optimize the project before starting the multi-state choreography epic.
The sprint therefore concentrated on the development environment and the
architecture: fixing the Docker images of the services, reducing their size,
restructuring the repository, configuring the robot connection through
environment variables, protecting the \texttt{master} branch and preparing a
CI pipeline that builds and publishes Docker images. Team documents are now
built from LaTeX and published to the team website automatically.
The team also analysed the causes of robot falls (SCRUM-54) and received
access to the faculty server (SCRUM-28).

During the sprint the team found inconsistencies between its methodologies:
the pull request approval rule is described differently in the Version
Control Methodology, the Technical Methodology and the Definition of Done,
and the use of the Testing status differs between the documents and
practice. These inconsistencies are being resolved in Sprint~2
(SCRUM-81 and SCRUM-83).

The tasks belonged to four epics: \emph{Development Environment Setup},
\emph{Architecture Improvement}, \emph{Project Management} and
\emph{Robot Stability and Fall Prevention}.


% ==================================================
% TASK OVERVIEW
% ==================================================

\section{Task Overview}

% Source: Jira sprint report of SCRUM Sprint 1.

\begin{naomitable}{Tasks of \sprintname{} (status as of \statusdate)}{tab:tasks}
    \footnotesize\setlength{\tabcolsep}{4pt}\renewcommand{\arraystretch}{1.25}
    \begin{tabularx}{\textwidth}{>{\raggedright\arraybackslash\mbox\bgroup}p{1.75cm}<{\egroup} Y L{2.9cm} L{1.75cm} L{0.5cm} L{1.15cm}}
        \rowcolor{naomiTableHeader}
        \textcolor{white}{\bfseries Key} &
        \textcolor{white}{\bfseries Task} &
        \textcolor{white}{\bfseries Assignee} &
        \textcolor{white}{\bfseries Status} &
        \textcolor{white}{\bfseries SP} &
        \textcolor{white}{\bfseries Sprint} \\
        SCRUM-27 & Update team website & A. Dmitriev & Done & 3 & 1 \\
        SCRUM-28 & Request and secure faculty server & T. Kuzin & Done & 3 & 0, 1 \\
        SCRUM-30 & Write Sprint 0 progress report and retrospective & R. Gainetdinov & Done & 3 & 1 \\
        SCRUM-36 & Automate LaTeX document build and publishing to team website & A. Zimnitski & Done & 5 & 1 \\
        SCRUM-43 & Install RobotSettings and Choregraphe & A. Dmitriev & Done & 2 & 1 \\
        SCRUM-51 & Fix missing dependency in voicecommandapi & D. Danilov & Done & 1 & 1 \\
        SCRUM-52 & Fix naorobotapi Docker image build & D. Danilov & Done & 2 & 1 \\
        SCRUM-54 & Analyze robot falls and existing fall detection in NAOqi & A. Akhmadinurov & Done & 5 & 1 \\
        SCRUM-55 & Set up CI pipeline to build and push Docker images & R. Gainetdinov & In Review & 5 & 1, 2 \\
        SCRUM-56 & Reduce Docker image sizes of microservices & A. Dmitriev & Done & 5 & 1 \\
        SCRUM-57 & Configure robot IP and port via environment variables & A. Dmitriev & Done & 2 & 1 \\
        SCRUM-58 & Simplify and document environment configuration & A. Zimnitski & In Review & 3 & 1, 2 \\
        SCRUM-59 & Protect master branch in naoMicroservices repository & A. Tikhoblazhenko & Done & 2 & 1 \\
        SCRUM-60 & Restructure naoMicroservices repository layout & D. Danilov & Done & 3 & 1 \\
        \arrayrulecolor{naomiTableRule}\bottomrule[1pt]
    \end{tabularx}
\end{naomitable}

\subsection{Scope Changes}

The sprint started on 26 October with four tasks (SCRUM-27, 28, 30, 36).
After the supervisor meeting and the analysis of the existing solution, the
team added SCRUM-43 on 29 October and SCRUM-51, 52, 54, 55, 56 and 57 on
30 October. On the same day all tasks were estimated in story points.
SCRUM-58, 59 and 60 were added on 2 November, after the decision to
refactor the project before the multi-state choreography epic.


% ==================================================
% TEAM PERFORMANCE
% ==================================================

\section{Team Performance}

\sprintname{} is the first sprint in which the team used story points.
The team completed \textbf{36 story points}.

\documentimage
    {velocity.png}
    {Velocity report after \sprintname{}}
    {fig:velocity}

The velocity report shows a commitment of 0 for \sprintname{}, because the
tasks were estimated only on 30 October, after the sprint had started.
Jira counts the commitment at the sprint start, so the real planned scope
(44 story points) is not visible in the chart. From Sprint~2 on, tasks are
estimated before the sprint starts. One sprint is not enough to determine
the team's velocity; it will become meaningful after several sprints.

\documentimage
    {burndown.png}
    {Sprint burndown chart of \sprintname{} (story points)}
    {fig:burndown}

For the same reason, the burndown chart stays at zero until 30 October and
then jumps to about 31 story points. The remaining work fell to 18 story
points by 1 November. On 2 November it rose to 26 story points when SCRUM-58,
59 and 60 were added, and it went back to 18 story points by 3 November.
The last tasks were closed on 8--9 November. The sprint ended with
8 story points remaining: SCRUM-55 and SCRUM-58, which were waiting for
pull request approval.


% ==================================================
% CUMULATIVE FLOW
% ==================================================

\section{Cumulative Flow}

The cumulative flow diagram covers \sprintname{} (26 October -- 9 November).
It counts all work items of the project, including epics.

\documentimage
    {cumulative-flow.png}
    {Cumulative flow diagram for \sprintname{}}
    {fig:cfd}

The green area (Done) grew from about 20 to about 43 work items during the
sprint. The purple area (To Do) grew on 30 October, 31 October and
2--3 November, when new tasks were created after the analysis of the
existing solution and the decision to refactor it. The large jump at the end
of the chart (9 November) corresponds to the planning of Sprint~2.
The bands for In Progress (cyan) and In Review (blue) stay narrow, so work
moved to Done without a bottleneck. At the end of the sprint the In Review
band widens slightly: these are the two tasks waiting for pull request
approval. The Testing status (yellow) was almost unused; how tasks without
code pass through Testing is being clarified in Sprint~2 (SCRUM-83).


\end{document}